\chapter{KAJIAN PUSTAKA}
\label{sec:chap2_kajian_pustaka}

\section{Hasil Penelitian Terdahulu}

Penelitian ini didukung oleh berbagai penelitian dan pengembangan sistem yang relevan dengan implementasi platform monitoring CCTV berbasis Internet of Things (IoT). Tinjauan terhadap penelitian terdahulu digunakan sebagai dasar untuk memahami perkembangan teknologi yang telah dilakukan sebelumnya, mengidentifikasi kelebihan dan keterbatasan sistem yang ada, serta menentukan kontribusi yang diberikan dalam penelitian ini.

\subsection{Pengembangan Aplikasi Mobile untuk Sistem Monitoring dan Pendeteksi Intruder Berbasis ESP32-CAM}

Penelitian yang dilakukan oleh Ali Akbar Alhabsyi mengembangkan aplikasi mobile berbasis Flutter untuk sistem monitoring dan pendeteksi intruder menggunakan kamera ESP32-CAM. Sistem yang dikembangkan memungkinkan pengguna melakukan pemantauan visual kamera secara langsung melalui perangkat smartphone, mengelola perangkat kamera, serta melakukan autentikasi pengguna dan pengelompokan kamera. Komunikasi data antara aplikasi dan perangkat dilakukan menggunakan protokol HTTP/TCP serta pertukaran data berformat JSON. Hasil penelitian menunjukkan bahwa aplikasi mampu menyediakan platform monitoring terpusat yang responsif dan mudah digunakan sehingga memudahkan pengguna dalam melakukan pengawasan dari berbagai lokasi. Namun, penelitian ini masih berfokus pada sisi aplikasi mobile dan belum membahas pemantauan telemetri perangkat secara mendalam maupun integrasi dashboard monitoring berbasis web secara real-time. \cite{ali2025}

\subsection{Pengembangan Platform MIOTVision Berbasis Laravel untuk Monitoring CCTV IoT}

Penelitian yang dilakukan oleh Ahmad Akmal Rijal mengembangkan platform MIOTVision berbasis web menggunakan framework Laravel untuk mengelola perangkat CCTV ESP32-CAM secara terpusat. Sistem ini mengintegrasikan broker MQTT EMQX, WebSocket Laravel Reverb, MySQL, dan MinIO untuk mendukung komunikasi data real-time, penyimpanan citra, serta pengelolaan metadata perangkat. Hasil penelitian menunjukkan bahwa sistem mampu menyediakan monitoring status perangkat secara real-time melalui mekanisme heartbeat, manajemen perangkat terpusat, pengelolaan log aktivitas, serta penyediaan pipeline data citra untuk kebutuhan pengembangan sistem yang lebih lanjut. Selain itu, sistem telah menerapkan mekanisme keamanan menggunakan Laravel Sanctum dan Role-Based Access Control (RBAC). Meskipun demikian, penelitian tersebut masih menemukan keterbatasan pada beberapa mekanisme sinkronisasi yang menggunakan HTTP Polling sehingga berpotensi menimbulkan overhead pada server dan latensi yang kurang konsisten ketika jumlah perangkat bertambah. Oleh karena itu, optimalisasi komunikasi berbasis WebSocket penuh dan pengembangan sistem pemantauan kondisi perangkat secara lebih komprehensif masih menjadi peluang pengembangan pada penelitian selanjutnya. \cite{rijal2025}

\subsection{Posisi Penelitian terhadap Penelitian Terdahulu}

Berdasarkan kedua penelitian terdahulu tersebut, telah tersedia aplikasi mobile untuk monitoring kamera ESP32-CAM serta platform web untuk pengelolaan perangkat CCTV berbasis IoT. Namun, masih terdapat beberapa aspek yang dapat dikembangkan lebih lanjut, yaitu integrasi website dan aplikasi mobile dalam satu platform yang terhubung secara real-time, implementasi komunikasi berbasis WebSocket secara penuh tanpa ketergantungan HTTP Polling, monitoring telemetri perangkat yang lebih komprehensif, implementasi notifikasi mobile otomatis, serta penerapan motion detection sebagai mekanisme deteksi aktivitas pada area pengawasan. Oleh karena itu, penelitian ini mengembangkan platform monitoring CCTV berbasis ESP32-CAM yang mengintegrasikan MQTT, WebSocket, motion detection, monitoring telemetri perangkat, notifikasi mobile, serta infrastruktur server terpusat berbasis Intel NUC untuk meningkatkan efektivitas pemantauan dan pengelolaan sistem secara keseluruhan.

\section{Dasar Teori}

\subsection{Internet of Things (IoT)}

Internet of Things (IoT) merupakan konsep yang memungkinkan berbagai perangkat fisik untuk saling terhubung melalui jaringan internet sehingga dapat melakukan pertukaran data secara otomatis tanpa memerlukan interaksi manusia secara langsung. Perangkat IoT umumnya dilengkapi dengan sensor, aktuator, dan modul komunikasi yang memungkinkan proses pengumpulan, pengiriman, serta pengolahan data secara real-time. Dalam penelitian ini, konsep IoT diterapkan pada sistem monitoring CCTV berbasis ESP32-CAM yang terhubung dengan server untuk mendukung proses pemantauan dan pengelolaan perangkat secara terpusat.

\subsection{ESP32-CAM}

ESP32-CAM merupakan modul mikrokontroler berbasis ESP32 yang telah dilengkapi dengan kamera OV2640 dan konektivitas Wi-Fi. Modul ini memiliki kemampuan untuk melakukan pengambilan citra, pemrosesan sederhana, serta pengiriman data melalui jaringan internet. Dengan ukuran yang relatif kecil dan biaya yang rendah, ESP32-CAM banyak digunakan pada berbagai aplikasi pengawasan dan monitoring berbasis IoT. Pada penelitian ini, ESP32-CAM berfungsi sebagai perangkat utama yang bertugas menangkap citra dan mengirimkannya ke server melalui protokol MQTT.

\subsection{MQTT (Message Queuing Telemetry Transport)}

MQTT merupakan protokol komunikasi ringan yang menggunakan arsitektur publish-subscribe. Protokol ini dirancang untuk perangkat dengan sumber daya terbatas dan jaringan dengan bandwidth rendah. Pada MQTT, komunikasi dilakukan melalui broker yang bertugas menerima pesan dari publisher dan meneruskannya kepada subscriber yang berlangganan pada topik tertentu. Dalam penelitian ini, MQTT digunakan sebagai media komunikasi utama antara perangkat ESP32-CAM dan server.

\subsection{EMQX Broker}

EMQX merupakan broker MQTT open-source yang dirancang untuk menangani komunikasi perangkat IoT dalam jumlah besar. Broker ini menyediakan berbagai fitur seperti autentikasi perangkat, manajemen koneksi, monitoring lalu lintas data, dan integrasi dengan berbagai sistem backend. Pada penelitian ini, EMQX digunakan sebagai pusat komunikasi MQTT yang menghubungkan perangkat ESP32-CAM dengan sistem server.

\subsection{WebSocket}

WebSocket merupakan protokol komunikasi yang memungkinkan pertukaran data dua arah secara real-time melalui satu koneksi yang bersifat persisten. Berbeda dengan mekanisme HTTP Polling yang memerlukan permintaan berulang dari klien, WebSocket memungkinkan server mengirimkan data secara langsung ketika terjadi perubahan informasi. Pada penelitian ini, WebSocket digunakan untuk memperbarui data monitoring pada website secara real-time tanpa perlu melakukan penyegaran halaman secara berulang.

\subsection{Laravel Framework}

Laravel merupakan framework berbasis PHP yang menerapkan pola arsitektur Model-View-Controller (MVC). Framework ini menyediakan berbagai fitur seperti routing, autentikasi, manajemen basis data, dan keamanan aplikasi yang memudahkan proses pengembangan perangkat lunak. Dalam penelitian ini, Laravel digunakan sebagai backend utama yang menangani proses pengelolaan data, komunikasi dengan broker MQTT, penyimpanan informasi perangkat, serta penyediaan API untuk website dan aplikasi mobile.

\subsection{React Native}

React Native merupakan framework pengembangan aplikasi mobile lintas platform yang memungkinkan pembuatan aplikasi Android dan iOS menggunakan bahasa pemrograman JavaScript. Framework ini memungkinkan penggunaan basis kode yang sama untuk berbagai platform sehingga proses pengembangan menjadi lebih efisien. Pada penelitian ini, React Native digunakan untuk membangun aplikasi mobile yang berfungsi menerima notifikasi serta memantau kondisi sistem CCTV secara real-time.

\subsection{MySQL Database}

MySQL merupakan sistem manajemen basis data relasional yang digunakan untuk menyimpan dan mengelola data secara terstruktur. MySQL mendukung penggunaan bahasa SQL untuk melakukan proses penyimpanan, pencarian, dan pengolahan data. Pada penelitian ini, MySQL digunakan untuk menyimpan informasi perangkat, data telemetri, metadata citra, data pengguna, serta informasi notifikasi yang dihasilkan oleh sistem.

\subsection{MinIO Object Storage}

MinIO merupakan sistem penyimpanan objek yang kompatibel dengan Amazon S3 API. MinIO dirancang untuk menyimpan data dalam jumlah besar dengan performa tinggi dan kemudahan integrasi dengan berbagai aplikasi. Pada penelitian ini, MinIO digunakan sebagai media penyimpanan citra yang dikirimkan oleh perangkat ESP32-CAM sehingga tidak membebani penyimpanan basis data.

\subsection{Docker Container}

Docker merupakan platform containerization yang memungkinkan aplikasi dijalankan dalam lingkungan terisolasi yang disebut container. Teknologi ini mempermudah proses deployment, pemeliharaan, dan pengelolaan layanan karena setiap komponen sistem dapat dijalankan secara independen. Dalam penelitian ini, Docker digunakan untuk menjalankan berbagai layanan seperti Laravel, MySQL, MinIO, dan EMQX pada server Intel NUC.

\subsection{Intel NUC Server}

Intel NUC merupakan komputer berukuran kecil yang memiliki performa cukup tinggi untuk menjalankan berbagai layanan server. Perangkat ini mendukung penggunaan sistem operasi Linux serta berbagai teknologi virtualisasi dan containerization. Pada penelitian ini, Intel NUC digunakan sebagai server utama yang menjalankan seluruh layanan backend dan penyimpanan data sistem monitoring CCTV.

\subsection{Cloudflared Tunnel}

Cloudflared Tunnel merupakan layanan yang memungkinkan akses ke server internal melalui jaringan Cloudflare tanpa perlu membuka port secara langsung pada router. Teknologi ini membantu meningkatkan keamanan sistem dengan menyembunyikan alamat IP server dari akses publik. Dalam penelitian ini, Cloudflared Tunnel digunakan untuk menyediakan akses website dan API secara aman melalui internet.

\subsection{Webmin}

Webmin merupakan aplikasi administrasi server berbasis web yang digunakan untuk mempermudah pengelolaan sistem operasi Linux. Melalui Webmin, administrator dapat melakukan konfigurasi layanan, pemantauan sumber daya, manajemen pengguna, dan berbagai tugas administrasi lainnya melalui antarmuka grafis. Pada penelitian ini, Webmin digunakan sebagai alat bantu dalam pengelolaan server Intel NUC.

\subsection{Motion Detection}

Motion Detection merupakan teknik yang digunakan untuk mendeteksi adanya perubahan pada suatu area pengamatan berdasarkan perbedaan antar frame citra. Metode ini memungkinkan sistem mengidentifikasi aktivitas atau pergerakan tanpa memerlukan proses klasifikasi objek yang kompleks. Pada penelitian ini, motion detection digunakan sebagai pemicu pengambilan gambar dan pengiriman notifikasi ketika terdeteksi aktivitas pada area pemantauan.

\subsection{Push Notification}

Push Notification merupakan mekanisme pengiriman pesan dari server ke perangkat pengguna secara langsung tanpa perlu pengguna membuka aplikasi terlebih dahulu. Teknologi ini banyak digunakan untuk memberikan informasi penting secara cepat dan real-time. Dalam penelitian ini, push notification digunakan untuk memberikan informasi kepada pengguna ketika sistem mendeteksi aktivitas atau kejadian tertentu pada area yang dipantau.

\subsection{Monitoring Telemetri Perangkat}

Monitoring telemetri perangkat merupakan proses pengumpulan dan pemantauan informasi operasional dari perangkat IoT secara berkala. Data telemetri digunakan untuk mengetahui kondisi perangkat sehingga administrator dapat melakukan pemeliharaan dan identifikasi masalah lebih cepat. Pada penelitian ini, data telemetri yang dipantau meliputi kekuatan sinyal jaringan, waktu aktif perangkat, penggunaan memori, penyebab reset perangkat, serta statistik keberhasilan proses pengambilan dan pengiriman data.

\subsubsection{Received Signal Strength Indicator (RSSI)}

RSSI merupakan parameter yang digunakan untuk menunjukkan kekuatan sinyal Wi-Fi yang diterima oleh perangkat. Nilai RSSI dapat digunakan untuk mengevaluasi kualitas koneksi jaringan dan membantu proses identifikasi gangguan komunikasi antara perangkat dan server.

\subsubsection{Uptime}

Uptime merupakan informasi yang menunjukkan lamanya perangkat beroperasi sejak terakhir kali dinyalakan atau mengalami reset. Parameter ini digunakan untuk mengetahui stabilitas operasional perangkat selama menjalankan sistem monitoring.

\subsubsection{Free Heap}

Free Heap merupakan jumlah memori yang masih tersedia pada perangkat untuk menjalankan proses dan menyimpan data sementara. Pemantauan Free Heap penting dilakukan untuk mendeteksi potensi kekurangan memori yang dapat menyebabkan kegagalan sistem.

\subsubsection{Reset Reason}

Reset Reason merupakan informasi yang menjelaskan penyebab perangkat melakukan proses restart. Informasi ini dapat digunakan untuk membantu proses diagnosis apabila terjadi gangguan atau ketidakstabilan pada perangkat.

\subsubsection{Capture Success Rate}

Capture Success Rate merupakan indikator yang menunjukkan tingkat keberhasilan perangkat dalam melakukan pengambilan citra. Nilai ini digunakan untuk mengevaluasi performa kamera dan memastikan proses akuisisi data berjalan dengan baik.

\subsubsection{Publish Success Rate}

Publish Success Rate merupakan indikator yang menunjukkan tingkat keberhasilan perangkat dalam mengirimkan data ke broker MQTT atau server. Parameter ini digunakan untuk menilai kualitas komunikasi data antara perangkat dan sistem backend.
